%! AP Statistics — Unit 1, Lesson 1.2 — Warm-Up KEY
\documentclass[10pt]{article}
\usepackage{apstats-article}
\usepackage{apstats-key}

%=============================================================================
\begin{document}

\pageheader{Unit 1, Lesson 1.2}{Warm-Up --- Answer Key}
\begin{tabularx}{\linewidth}{X X X}
Name:\;\ans{Key} & Date:\; & Period:\;
\end{tabularx}

\vspace{0.28in}

{\large\bfseries\color{navy} Part 1 \quad Revisiting Lesson 1.1}

\vspace{0.12in}

\noindent The table below shows data from a commute study of five workers.

\vspace{0.1in}

\renewcommand{\arraystretch}{1.6}
\begin{tabularx}{\linewidth}{@{} >{\bfseries}p{0.18\linewidth}
                                    C{0.16\linewidth} C{0.18\linewidth}
                                    C{0.22\linewidth} C{0.16\linewidth} @{}}
\rowcolor{sky}
\textbf{Name} & \textbf{Mode} & \textbf{Commute (min)} &
\textbf{Satisfaction (1--5)} & \textbf{Days/wk} \\
\midrule
Aisha   & Bus    & 42 & 3 & 5 \\
Ben     & Car    & 18 & 5 & 5 \\
Carmen  & Train  & 61 & 2 & 4 \\
DeShawn & Bike   & 25 & 4 & 3 \\
Elena   & Car    & 34 & 4 & 5 \\
\end{tabularx}

\vspace{0.2in}

\begin{enumerate}

  \item Which column(s) record words or labels rather than numbers?\\[6pt]
        \ans{Mode (Bus, Car, Train, Bike) --- these are category labels, not numbers.
        Name is also a label, but it is not really a variable of interest.}

  \item Which column(s) record numbers you could reasonably average?
        Find the average for one of them.\\[6pt]
        \ans{Commute time and Days/week. For example, average commute time $= (42+18+61+25+34)/5 = 180/5 = 36$ minutes.
        Average days/week $= (5+5+4+3+5)/5 = 22/5 = 4.4$ days per week.}

  \item Satisfaction is recorded as a 1--5 scale. Can you average those values?
        Does the average \textit{mean} anything useful in this context?\\[6pt]
        \ans{You can compute the average: $(3+5+2+4+4)/5 = 18/5 = 3.6$.
        Whether it is \emph{meaningful} is debatable. The numbers label ordered levels of
        satisfaction, but the gap between ``1'' and ``2'' may not equal the gap between
        ``4'' and ``5.'' Many statisticians treat this as categorical (ordinal); others
        treat it as quantitative for convenience. Key takeaway: context and judgment matter.}

  \item Sort the four variables:

\vspace{0.1in}
\renewcommand{\arraystretch}{2.2}
\begin{tabularx}{\linewidth}{@{} |>{\centering\arraybackslash}X
                                    |>{\centering\arraybackslash}X| @{}}
\hline
\rowcolor{sky}
\textbf{``Word / label'' variables} & \textbf{``Number'' variables} \\
\hline
\ans{Mode of transport} & \ans{Commute time (min), Days/week,} \\
& \ans{Satisfaction (debatable)} \\
\hline
\end{tabularx}

\end{enumerate}

\begin{teachernote}
\textbf{Discussion point:} The Satisfaction variable is intentionally ambiguous ---
it is the entry point for the lesson's key idea that some numbers are not truly
quantitative. Accept either answer for Satisfaction as long as the student gives a
reason. If most students put it in the ``number'' column without comment, use that
as the hook to open the lesson.
\end{teachernote}

\vspace{0.2in}

{\large\bfseries\color{navy} Part 2 \quad Connecting to Today}

\vspace{0.12in}

\begin{tcolorbox}[colback=greenbg, colframe=greenacc, boxrule=0.8pt, arc=2mm,
    left=4mm, right=4mm, top=3mm, bottom=3mm]
\small
\begin{itemize}[leftmargin=1.3em, itemsep=8pt]
  \item The two groups you just created have names in statistics:
        \textbf{categorical} and \textbf{quantitative}. Today we formalize this.
  \item Variable type determines \emph{which displays and summaries make sense}.
        Lessons 1.3--1.4 focus on categorical; Lessons 1.5--1.9 focus on quantitative.
  \item Identifying variable type is the \emph{first} question on every AP free-response
        problem that involves data.
\end{itemize}
\end{tcolorbox}

\vspace{0.16in}

\noindent\textcolor{slate}{\small One thing you're not sure about yet, or one question going into today's lesson:}\\[6pt]
\ans{Answers vary --- look for genuine uncertainty. Common examples: ``What makes something truly quantitative?'' / ``Is age categorical or quantitative?'' / ``What is a distribution?'' These are all productive questions to validate and connect to today's content.}

\end{document}
